\documentclass[11pt]{article}

\usepackage[T1]{fontenc}
\usepackage{lmodern}
\usepackage[margin=1in]{geometry}
\usepackage{amsmath,amssymb,amsthm}
\usepackage{booktabs}
\usepackage{microtype}
\usepackage[hidelinks]{hyperref}
\hypersetup{
  pdftitle={Even-Parameter Periods for Reversible Elementary Second-Order Cellular Automaton 115: A Computer-Assisted Proof},
  pdfauthor={Xiaojian Zeng},
  pdfsubject={Exact least periods for RESOCA rule 115},
  pdfkeywords={cellular automata, reversible dynamics, exact period, finite-state certificate}
}

\newtheorem{theorem}{Theorem}
\newtheorem{lemma}[theorem]{Lemma}
\newtheorem{proposition}[theorem]{Proposition}
\newtheorem{corollary}[theorem]{Corollary}
\theoremstyle{remark}
\newtheorem{remark}[theorem]{Remark}

\newcommand{\Ftwo}{\mathbb F_2}
\newcommand{\InnerSet}{\mathcal I}
\newcommand{\OuterSet}{\mathcal O}
\newcommand{\GenericSet}{\mathcal G}
\newcommand{\NextSet}{\mathcal N}

\title{Even-Parameter Periods for Reversible Elementary\\
Second-Order Cellular Automaton 115:\\A Computer-Assisted Proof}
\author{Xiaojian Zeng\\
\small PBC School of Finance, Tsinghua University, 100190\\
\small \texttt{xiaojianzeng@foxmail.com}}
\date{}

\begin{document}
\maketitle

\begin{abstract}
Formenti and Kamilya conjectured an exact formula for the least period of the
finite initial-configuration family denoted by $1(01)^k$ under reversible
elementary second-order cellular automaton rule 115.  We prove their
even-parameter formula for every $i\geq 0$:
\[
 P(2i)=\frac{2\,4^{2i+2}-7\,4^{i+1}+5}{3}.
\]
The proof samples every third time step and encodes pairs of adjacent sites by
base-four blocks.  An exact finite-window certificate splits the sampled orbit
into two independent finite path-language permutations.  A second exact
finite-state certificate proves that each permutation is the cyclic successor
in an explicitly defined colexicographic order.  The resulting factor periods
are coprime, and the period-three background supplies the final leastness
argument.  The computational part uses only exhaustive Boolean/integer checks;
its predicates and all-length induction are described explicitly.
\end{abstract}

\section{Statement and context}

Elementary cellular automata use a binary state set, a one-dimensional lattice,
and a radius-one local rule; we use the standard Wolfram numbering convention
\cite{Wolfram1983}.  An early explicit second-order reversible construction
appears in Margolus's work on physics-like models of computation
\cite[Appendix~A]{Margolus1984}.  Toffoli, Capobianco, and Mentrasti later
gave an inversion scheme by representing any second-order cellular automaton
as a functionally isomorphic lattice gas \cite{ToffoliCapobiancoMentrasti2004}.
The broader reversible-cellular-automaton setting is surveyed by Toffoli and
Margolus \cite{ToffoliMargolus1990}, while modern work on higher-order cellular
automata includes decidability results for the linear subclass
\cite{DennunzioEtAl2019}.  The elementary RESOCA subclass used here was
introduced and studied systematically by Formenti and Kamilya
\cite{FormentiKamilya2024}.

Let $X=\{0,1\}^{\mathbb Z}$, with addition denoted by $\oplus$.  Wolfram rule
115 is the map $\delta_{115}:\{0,1\}^3\to\{0,1\}$ whose values, in the
neighborhood order
\[
111,110,101,100,011,010,001,000,
\]
are respectively
\[
0,1,1,1,0,0,1,1.
\]
The associated reversible elementary second-order cellular automaton is
\begin{equation}\label{eq:F}
 F(b,c)=(c,d),\qquad
 d_j=b_j\oplus\delta_{115}(c_{j-1},c_j,c_{j+1})
 \quad(j\in\mathbb Z).
\end{equation}
Its inverse is immediate: if $F(b,c)=(c,d)$, then
\begin{equation}\label{eq:inverse}
 b_j=d_j\oplus\delta_{115}(c_{j-1},c_j,c_{j+1}).
\end{equation}

Let $u\in X$ have $u_0=1$ and $u_j=0$ for $j\neq0$.  For $k\geq0$, let
$c^{(k)}$ be the centered configuration whose finite word is $1(01)^k$;
equivalently,
\[
 c^{(k)}_j=1 \quad\Longleftrightarrow\quad
 j\in\{-k,-k+2,\ldots,k\}.
\]
Write $P(k)$ for the least positive period of $(u,c^{(k)})$ under $F$.
Formenti and Kamilya proved that every finite initial pair is periodic and that
its least period is a multiple of three.  They then stated the formula below as
Conjecture~3 \cite[Sec.~6.3]{FormentiKamilya2026}.

\begin{theorem}\label{thm:main}
For every integer $i\geq0$, the least positive period of
$(u,c^{(2i)})$ is
\[
 P(2i)=\frac{2\,4^{2i+2}-7\,4^{i+1}+5}{3}.
\]
In particular, the formula includes $P(0)=3$.
\end{theorem}

Our contribution is an all-$i$ equality and a proof of leastness.  The odd
subfamily conjecture in the source paper is not addressed here.  A separate
recent preprint reports a symbolic proof of that neighboring odd-parameter
formula using a finite colored-tree walk \cite{Zeng2026SSRN7379379}; the
factor/colex certificate used below is different.

\section{Three-step blocks and path languages}

Put $H=F^3$.  Under $F$, the uniform site-pair follows
\begin{equation}\label{eq:background}
 (0,0)\to(0,1)\to(1,0)\to(0,0).
\end{equation}
Thus times divisible by three again have zero background.  Encode a site-pair
$(b_j,c_j)$ as $2b_j+c_j\in\{0,1,2,3\}$, pair the adjacent sites
$(2r-1,2r)$, and encode the resulting two symbols by
\begin{equation}\label{eq:blockcode}
 [r]=4(2b_{2r-1}+c_{2r-1})+(2b_{2r}+c_{2r})\in\{0,\ldots,15\}.
\end{equation}
Three radius-one updates make the output block $[r]$ depend on the five input
blocks $[r-2],\ldots,[r+2]$.

We now define the finite path languages that describe the sampled orbit.  Let
\[
 \GenericSet=\{1,2,4,6,8,9,12,13,14,15\},
\]
and define allowed successors by
\begin{center}
\begin{tabular}{c@{\qquad}l}
$a$ & $\NextSet(a)$\\
$1,9,13$ & $\{1,2,8,9\}$\\
$2,6,14$ & $\{1,2,4,6\}$\\
$4,8,12$ & $\{4,6,8,9\}$\\
$15$ & $\GenericSet$\\
$0$ & $\{4,6,8,9\}$\\
$3$ & $\GenericSet$\\
\end{tabular}
\end{center}
Also set
\[
 \OuterSet=\{1,2,12,13,14,15\},\qquad
 \InnerSet=\{1,2,6,9,13,14,15\}.
\]
For $i\geq1$, let $L_i$ contain the words $x_1\cdots x_i$ over $\GenericSet$
such that $x_1\in\OuterSet$, $x_i\in\InnerSet$, and
$x_{j+1}\in\NextSet(x_j)$.  Let $L_0=\{\epsilon\}$.
For $i\geq1$, let $R_i$ contain the words $x_0x_1\cdots x_i$ such that
$x_0\in\{0,3\}$, $x_1,\ldots,x_i\in\GenericSet$,
$x_{j+1}\in\NextSet(x_j)$, and $x_i\in\InnerSet$.  Finally set
$R_0=\{(3)\}$.

Under the encoding \eqref{eq:blockcode}, the initial sampled state for family
parameter $k=2i$ splits at the central interface into $i$ left blocks and
$i+1$ right blocks as
\begin{equation}\label{eq:initial-factors}
 (1,\ldots,1)\;(3,1,\ldots,1)
 \in L_iR_i.
\end{equation}
Indeed, the left blocks cover the sites $[-2i-1,-2]$ and the right blocks
cover $[-1,2i]$.

\begin{lemma}[Path counts]\label{lem:counts}
For every $i\geq0$,
\begin{equation}\label{eq:AB}
 |L_i|=A_i=\frac{4^{i+1}-1}{3},\qquad
 |R_i|=B_i=\frac{2\,4^{i+1}-5}{3}.
\end{equation}
\end{lemma}

\begin{proof}
Index a $10\times10$ zero-one matrix $M$ by the elements of $\GenericSet$ and
put $M_{ab}=1$ exactly when $b\in\NextSet(a)$.  Let $e$ be the column
indicator of $\InnerSet$.  Direct row substitution using the displayed successor
table gives
\begin{equation}\label{eq:minpoly}
 (M^2-5M+4I)e=0.
\end{equation}
Consequently every path-count sequence obtained by applying a fixed starting
row vector to $M^ne$ satisfies
$C_{n+2}=5C_{n+1}-4C_n$.  Directly from the same table,
\[
 |L_0|,|L_1|,|L_2|=1,5,21,
 \qquad
 |R_0|,|R_1|,|R_2|=1,9,41.
\]
The two expressions in \eqref{eq:AB} have these initial values and satisfy the
recurrence, proving the claim.
\end{proof}

\section{Exact invariance and factorization}

Finite-state procedures that derive global properties of one-dimensional
cellular automata from their local maps are classical; the injectivity and
surjectivity decision procedures of Amoroso and Patt are an early precedent
\cite{AmorosoPatt1972}.  The certificate below is not an application of those
procedures: it is specialized to the two invariant path languages defined
above and to their zero-padded finite boundaries.

The next lemma is computer-assisted.  We state both the mathematical claim and
the finite verification that implies it.

\begin{lemma}[Block factorization]\label{lem:factorization}
For every $i\geq0$, applying the zero-padded block map $H$ sends $L_i$ to
$L_i$ and $R_i$ to $R_i$, with zero output outside the corresponding block
interval.  Moreover,
\begin{equation}\label{eq:decouple}
 H(lr)=H(l)H(r)\qquad(l\in L_i,\ r\in R_i),
\end{equation}
where the two factors on the right are evolved with zero padding.  Each
induced factor map is a permutation.
\end{lemma}

\begin{proof}
Define the five-block local map $h:\{0,\ldots,15\}^5\to
\{0,\ldots,15\}$ by decoding the five blocks into ten site-pairs,
applying \eqref{eq:F} exactly three times, and returning the central two-site
block.  An exhaustive program evaluates all $16^5=1{,}048{,}576$ inputs.

Because the input languages are one-step path languages and $h$ has block
radius two, every possible interior output symbol occurs in an admissible
five-block input path, and every possible adjacent pair of outputs occurs in
an admissible six-block input path.  The certificate checks all $3{,}070$
five-block and $12{,}286$ six-block generic paths.  These checks prove that
interior outputs belong to $\GenericSet$ and retain the successor relation.

Only the first two and last two output blocks can meet a padded boundary.  The
certificate exhausts $8{,}190$ left prefixes, $4{,}094$ right prefixes, and
$8{,}191$ generic terminal suffixes of length six.  It separately tests all
shorter words for which the two boundary windows overlap: $1{,}818$ left
words and $3{,}630$ right words.  It checks both the relevant initial or
terminal condition and that all four possible spill blocks are zero.  This
proves all-length invariance.

At the central interface, block radius two means that only the last two left
outputs and the first two right outputs can depend on the other factor.  All
$511$ generic left terminal contexts and $254$ right initial contexts of
length four are combined, giving $129{,}794$ interface comparisons with the
corresponding zero-padded factor outputs.  A further $15{,}272$ shorter factor
pairs cover cases where an outer boundary also lies in the local window.  All
comparisons give \eqref{eq:decouple}.

These are exhaustive local implications, not tests through a bounded word
length.  Finally, \eqref{eq:inverse} makes $H$ injective.  Its restrictions are
injective self-maps of finite sets, hence permutations.
\end{proof}

Lemma~\ref{lem:factorization} does not by itself show that either permutation
is a single cycle.  That is the purpose of the explicit order below.

\section{Colexicographic successor cycles}

Adjoin a formal right endpoint $0$.  For each possible right-hand symbol $t$,
order its generic predecessors as follows:
\begin{align*}
O_0&=(1,9,13,14,6,2,15),\\
O_1=O_2&=(14,6,2,15,1,9,13),\\
O_4=O_6&=(8,12,4,14,6,2,15),\\
O_8=O_9&=(15,1,9,13,8,12,4),\\
O_{12}=O_{13}=O_{14}=O_{15}&=(15).
\end{align*}
For $t=0$, this is the terminal set $\InnerSet$ in its chosen order.  In every
other case, comparison with the successor table shows that $O_t$ contains
exactly those $a\in\GenericSet$ for which $t\in\NextSet(a)$.

For $X\in\{L,R\}$, let $E_m^X(t)$ be an ordered list of length-$m$ prefixes.
For the left lists, $E_1^L(t)$ consists of the one-letter words $(a)$ with
$a\in O_t\cap\OuterSet$, in inherited order.  For the right lists, define
\[
 C_8=C_9=(3,0),\qquad C_4=C_6=(0,3),\qquad C_t=(3)
 \quad\text{otherwise},
\]
for generic $t$, put $E_1^R(t)=((a):a\in C_t)$, and set
$E_1^R(0)=((3))$.  For $m>1$ on either side, define recursively
\begin{equation}\label{eq:Erec}
 E_m^X(t)=\mathop{\operatorname{concat}}_{a\in O_t}
 \{(v,a):v\in E_{m-1}^X(a)\},
\end{equation}
where $\operatorname{concat}$ denotes concatenation in the displayed order.

Since every predecessor occurs exactly once in $O_t$, induction on $m$ shows
that these lists contain every corresponding admissible path exactly once.
In particular,
\begin{equation}\label{eq:E-language}
 E_i^L(0)=L_i\quad(i\geq1),\qquad
 E_{i+1}^R(0)=R_i\quad(i\geq1),
\end{equation}
with the already specified exceptional $i=0$ factors.

\begin{lemma}[Cyclic successor]\label{lem:successor}
For every admissible length on each side, the zero-padded map $H$ sends each
word of $E_m^X(0)$ to the next word of that list, cyclically.
\end{lemma}

\begin{proof}
We give the exact finite-state induction used to certify the claim.  Let
$Q_m^X(t)$ be the set of aligned pairs of adjacent words in $E_m^X(t)$,
excluding wraparound.  From \eqref{eq:Erec}, it obeys the exact recurrence
\begin{equation}\label{eq:Qrec}
\begin{split}
Q_m^X(t)={}&\bigcup_{a\in O_t}
 \{((v,a),(w,a)):(v,w)\in Q_{m-1}^X(a)\}\\
&{}\cup\bigcup_{a,b\text{ consecutive in }O_t}
 \{((\operatorname{last}E_{m-1}^X(a),a),
    (\operatorname{first}E_{m-1}^X(b),b))\}.
\end{split}
\end{equation}
The first and last words themselves satisfy the one-symbol recursion obtained
from the first and last elements of $O_t$.

For a collection of equal-length aligned word pairs, retain all pair factors
of lengths at most five, together with their five-pair prefixes and suffixes
(or the whole word if shorter).  Appending a fixed pair of symbols updates
this summary exactly: old factors remain, new factors are precisely the
suffixes of length at most four followed by the appended pair, prefixes stay
fixed once the length is at least five, and suffixes shift by one.  Union and
the bridge pairs in \eqref{eq:Qrec} therefore give an exact deterministic
summary recurrence.

The induction state contains these summaries for every $Q_m^X(t)$, the
first and last words truncated to their first and last five symbols, and the
same summaries for every aligned pair of such extrema.  These latter records
determine every bridge and the cyclic last-to-first pair at the next length.
Thus the retained state contains all and only the finite information needed
to construct the next state and to test a radius-two block map.

For each aligned transition $x\to y$ in the summary of $E_m^X(0)$, including
the cyclic transition from the last word to the first, the certificate checks
\[
 h(x_{j-2},x_{j-1},x_j,x_{j+1},x_{j+2})=y_j
\]
at every interior position and checks the two output and two spill records at
each zero-padded boundary.  Lengths below five are checked directly.  The
complete local-state signatures at left lengths $7$ and $8$ agree, and those
at right lengths $8$ and $9$ agree.  The program checks $44{,}301$ left and
$48{,}650$ right summary records before these repetitions.  Since the summary
update is deterministic and the repetitions occur beyond the five-block
window length, induction gives the same verified signature, hence all the
same local transition identities, at every later length.

Every earlier length was checked directly or by the exact summary.  Therefore
$H(x)=y$ for each adjacent pair in every list, including wraparound, which is
the asserted cyclic-successor action.
\end{proof}

\begin{corollary}\label{cor:factorperiods}
The least $H$-periods of the two factors in \eqref{eq:initial-factors} are
$A_i$ and $B_i$, respectively.
\end{corollary}

\begin{proof}
By \eqref{eq:E-language}, each ordered list contains every word of its path
language exactly once.  Lemma~\ref{lem:successor} makes its induced permutation
one cycle.  Its cycle length is therefore its cardinality, given by
Lemma~\ref{lem:counts}.  The singleton factors at $i=0$ have period one under
$H$ and obey the same formulas.
\end{proof}

\section{Proof of Theorem~\ref{thm:main}}

Lemma~\ref{lem:factorization} makes the two factor evolutions independent.
Corollary~\ref{cor:factorperiods} gives their least sampled periods $A_i$ and
$B_i$.  Since
\begin{equation}\label{eq:coprime}
 B_i=2A_i-1,
\end{equation}
we have $\gcd(A_i,B_i)=1$.  Hence the concatenated sampled state has least
$H$-period $A_iB_i$.

It remains to pass from $H=F^3$ back to $F$.  Equation
\eqref{eq:background} shows that infinitely many sites outside the finite
defect are in distinct uniform phases at times congruent to $0,1,2$ modulo
three.  Therefore the initial pair cannot return at a time not divisible by
three.  Its least $F$-period is consequently
\begin{align*}
3A_iB_i
&=3\left(\frac{4^{i+1}-1}{3}\right)
       \left(\frac{2\,4^{i+1}-5}{3}\right)\\
&=\frac{2\,4^{2i+2}-7\,4^{i+1}+5}{3}.
\end{align*}
For $i=0$, both sampled factors are fixed and the genuine background cycle
has length three, so the least period is exactly $3$.  This proves the theorem.

\section{Reproducibility and assistance disclosure}

The block-factorization certificate and colex-successor certificate are the
plain Python 3 programs
\begin{center}
\texttt{evidence/research\_block\_transducer.py},\qquad
\texttt{evidence/research\_colex\_certificate.py}.
\end{center}
They use exact Boolean and integer operations and Python's standard library;
neither floating-point arithmetic nor randomized testing is used.  From the
project root they can be rerun with
\begin{verbatim}
python3 evidence/research_block_transducer.py
python3 evidence/research_colex_certificate.py
\end{verbatim}
The first program aborts on a failed local-rule, boundary, language, or
interface assertion.  The second implements \eqref{eq:Qrec}, the five-local
summary update, all boundary/spill checks, and the fixed-signature induction
described in Lemma~\ref{lem:successor}.  Their supplied transcripts record a
passing run and the counts reported above.  These computations certify only
the precisely stated finite implications; the path-count, coprimality, and
leastness arguments are mathematical deductions given in the text.

AI assistance was used in proof search, construction and review of the exact
certificate programs, and preparation of this exposition.  The certificate
outputs are reproducible artifacts, not substitutes for the stated logical
reductions, and an AI review is not represented as independent peer review.

\bibliographystyle{plain}
\bibliography{references}

\end{document}
